\chapter{Notation, Assumptions, And Reading Rules}
\label{chap:notation}

\section{Purpose Of These Notes}
\label{sec:notation-purpose}

These notes are written as lecture notes rather than as a compressed reference.
Each chapter states its assumptions, defines its symbols before use whenever
possible, and derives the central equations from a clearly named structure:
action principles, symplectic geometry, symmetry, or continuum balance laws.

The guiding convention is simple: when an equation appears, the reader should
know what space it lives on, which variables are independent, which assumptions
make it valid, and what physical information it carries.

The notes are meant to be self-contained at the level of a first serious pass
through advanced mechanics. When a result is used later, the earlier occurrence
should have supplied the model, the variables, the assumptions, and the
derivation strategy. The reader should never have to guess whether an equation
is a definition, an equation of motion, a constraint, a conservation law, or a
coordinate representation of a geometric object.

Self-contained does not mean that every calculation starts from elementary
calculus. It means that every new structure is introduced with enough context
to be used responsibly. For example, when the notes invoke a symplectic form,
they state the phase space and the sign convention; when they invoke a stress
tensor, they state the balance law and the constitutive assumption; when they
invoke a connection, they say which frame choice is being compared.

The recurring reading questions are:
\[
\begin{gathered}
  \text{What is the configuration or phase space?}\\
  \text{What structure is being preserved or varied?}\\
  \text{Which assumptions make the calculation legal?}\\
  \text{What physical or geometric information survives a change of coordinates?}
\end{gathered}
\]
These questions are deliberately repeated in different contexts. They are the
thread connecting particles, rigid bodies, celestial mechanics, fluids, and
deforming continua.

\section{Exposition Standard Used In The Notes}
\label{sec:exposition-standard}

Each major derivation is written to answer five questions in order:
\[
\begin{array}{rcl}
\text{model} &:& \text{what physical idealization is being made?}\\
\text{space} &:& \text{where do the variables live?}\\
\text{principle} &:& \text{what variational, Hamiltonian, or balance statement is used?}\\
\text{calculation} &:& \text{which algebraic or analytic step produces the equation?}\\
\text{interpretation} &:& \text{what does the result say invariantly?}
\end{array}
\]
When a derivation feels long, this list is the intended guide. A displayed
equation should not be read as an isolated fact; it is either part of the
model, a definition of notation, a consequence of a principle, or an
interpretation of the result. The surrounding prose is meant to identify which
role the equation is playing.

The notes also distinguish three kinds of equality. A definition introduces a
symbol, as in \(p_i=\partial L/\partial\dot q^i\). An identity holds because of
previous definitions, as in \(\omega=-\dd\Theta\). An equation of motion holds
only along physical trajectories, as in \(\dot p_i=-\partial H/\partial q^i\).
Many confusions in mechanics come from silently moving an equation from one
category to another.

\section{How Notation Is Controlled}
\label{sec:notation-control}

Notation in mechanics is necessarily overloaded: \(I\) can mean a time
interval, an inertia tensor, or an action variable; \(\omega\) can mean a
symplectic form, an angular velocity, a frequency vector, or vorticity. The
notes therefore use the following discipline.

\begin{enumerate}
\item A symbol has no permanent meaning unless it appears in the global ledger
below or in the notation table at the beginning of a chapter.
\item If a standard symbol is reused with a different meaning, the local
definition overrides the global convention only inside that section or chapter.
\item Bold or vector notation indicates Euclidean vectors only when the ambient
space and metric have been named.
\item Upper/lower indices in finite-dimensional mechanics indicate coordinate
indices on a manifold and its cotangent bundle. In continuum chapters, capital
indices such as \(A,B\) label material coordinates and lower-case indices such
as \(i,j\) label spatial coordinates.
\item Every theorem, derivation, example, and demo should be readable together
with the nearest notation table and the assumptions stated immediately before
the calculation.
\end{enumerate}

\section{Standing Mathematical Assumptions}
\label{sec:standing-assumptions}

\begin{assumption}[Smoothness]
Unless explicitly stated otherwise, all functions, curves, vector fields, and
maps are assumed smooth enough for the differentiations and integrations being
performed. For most finite-dimensional chapters, \(C^2\) regularity is enough.
For continuum chapters, additional differentiability is assumed locally.
\end{assumption}

\begin{assumption}[Finite-dimensional mechanics]
When the configuration space is a finite-dimensional manifold \(Q\), a motion is
a curve \(q:I\to Q\) on a time interval \(I=[t_0,t_1]\). The velocity
\(\dot q(t)\) lies in the tangent space \(T_{q(t)}Q\). Local coordinates are
written \(q^1,\ldots,q^n\).
\end{assumption}

\begin{assumption}[Endpoint variations]
For Hamilton's principle in its basic form, variations of the path satisfy
\(\delta q(t_0)=\delta q(t_1)=0\). Free-endpoint and boundary-force variants are
important, but they are not the default convention.
\end{assumption}

\begin{assumption}[Regular Lagrangians]
The Legendre transform from velocities to momenta is assumed locally invertible
whenever we pass from \(L(q,\dot q,t)\) to \(H(q,p,t)\). In coordinates this
means the matrix
\[
  \frac{\partial^2L}{\partial \dot q^i\partial \dot q^j}
\]
is nonsingular. Singular Lagrangians lead to constraints and gauge theories in
the Dirac sense; those require a separate treatment.
\end{assumption}

\begin{assumption}[Einstein summation]
Repeated upper/lower coordinate indices are summed unless stated otherwise:
\[
  p_i\dot q^i=\sum_{i=1}^{n}p_i\dot q^i.
\]
In Euclidean vector formulas, dots and cross products use the standard metric
and orientation.
\end{assumption}

\section{Global Notation Ledger}
\label{sec:global-notation-ledger}

This table is a reference ledger, not a replacement for local definitions. The
``first reference'' column points to the chapter where the symbol family is
used most systematically.

\small
\begin{longtable}{p{0.18\textwidth}p{0.52\textwidth}p{0.22\textwidth}}
\toprule
Symbol & Meaning and assumptions & First reference\\
\midrule
\endhead
\bottomrule
\endfoot
\(t\) & time variable; dots denote \(\dd/\dd t\) unless a section says otherwise & all chapters\\
\(I=[t_0,t_1]\) & time interval for a path or field evolution & Chapter~\ref{chap:lag-ham}\\
\(Q\) & finite-dimensional configuration manifold & Chapter~\ref{chap:lag-ham}\\
\(q(t)\) & configuration path \(I\to Q\) & Chapter~\ref{chap:lag-ham}\\
\(q^i\) & local coordinates on \(Q\), \(i=1,\ldots,n\) & Chapter~\ref{chap:lag-ham}\\
\(\dot q^i\) & generalized velocities \(\dd q^i/\dd t\) & Chapter~\ref{chap:lag-ham}\\
\(TQ\) & tangent bundle, the space of \((q,\dot q)\) & Chapter~\ref{chap:lag-ham}\\
\(T^*Q\) & cotangent bundle, the phase space of \((q,p)\) & Chapter~\ref{chap:lag-ham}\\
\(L\) & Lagrangian, usually \(T-V\) but not assumed to be natural unless stated & Chapter~\ref{chap:lag-ham}\\
\(S[q]\) & configuration-space action functional \(\int L\,\dd t\) & Chapter~\ref{chap:lag-ham}\\
\(\eta\) & variation vector field \(\delta q\) along a path & Chapter~\ref{chap:lag-ham}\\
\(p_i\) & canonical momentum \(\partial L/\partial \dot q^i\) & Chapter~\ref{chap:lag-ham}\\
\(H\) & Hamiltonian on phase space; in autonomous conservative examples it is energy & Chapter~\ref{chap:lag-ham}\\
\(\Theta\) & canonical one-form on \(T^*Q\), \(\Theta=p_i\,\dd q^i\) & Chapter~\ref{chap:lag-ham}\\
\(\omega\) & symplectic form in Hamiltonian chapters; other uses are local and stated & Chapter~\ref{chap:lag-ham}\\
\(\{f,g\}\) & Poisson bracket of observables \(f,g\) & Chapter~\ref{chap:lag-ham}\\
\(X_H\) & Hamiltonian vector field determined by \(\iota_{X_H}\omega=\dd H\) & Chapter~\ref{chap:lag-ham}\\
\(G\) & Lie group in symmetry sections; gravitational constant in celestial sections & Chapters~\ref{chap:lag-ham}, \ref{chap:perturbation-chaos}\\
\(\mathfrak g\) & Lie algebra of a Lie group \(G\) & Chapters~\ref{chap:lag-ham}, \ref{chap:elastic-gauge}\\
\(\xi_Q\) & infinitesimal generator on \(Q\) associated with \(\xi\in\mathfrak g\) & Chapter~\ref{chap:lag-ham}\\
\(J_\xi\) & Noether conserved quantity associated with \(\xi\) & Chapter~\ref{chap:lag-ham}\\
\(R\) & rotation matrix in rigid-body chapters; radius or gas constant only when locally declared & Chapter~\ref{chap:rigid-integrability}\\
\(\SO(3)\) & group of orientation-preserving orthogonal \(3\times3\) matrices & Chapter~\ref{chap:rigid-integrability}\\
\(\Omega\) & body angular velocity vector; not the same as symplectic \(\omega\) & Chapter~\ref{chap:rigid-integrability}\\
\(\widehat\Omega\) & skew matrix satisfying \(\widehat\Omega v=\Omega\times v\) & Chapter~\ref{chap:rigid-integrability}\\
\(\omega_{\mathrm{space}}\) & spatial angular velocity when needed to avoid conflict with symplectic \(\omega\) & Chapter~\ref{chap:rigid-integrability}\\
\(I\) & inertia tensor or action variable depending on context; local table fixes meaning & Chapters~\ref{chap:rigid-integrability}, \ref{chap:perturbation-chaos}\\
\(M\) & body angular momentum in rigid-body chapters & Chapter~\ref{chap:rigid-integrability}\\
\(J\) & spatial angular momentum in rigid-body/celestial chapters; Jacobian determinant in continuum chapters & Chapters~\ref{chap:rigid-integrability}, \ref{chap:fluid}\\
\(F_i\) & commuting first integrals in Liouville theory; deformation gradient \(F^i{}_A\) in continua & Chapters~\ref{chap:rigid-integrability}, \ref{chap:elastic-gauge}\\
\(\gamma_j\) & homology cycle on an invariant torus; material loop in fluid contexts & Chapters~\ref{chap:rigid-integrability}, \ref{chap:fluid}\\
\(I_j\) & action variable \((2\pi)^{-1}\oint_{\gamma_j}p_i\,\dd q^i\) & Chapter~\ref{chap:rigid-integrability}\\
\(\theta^j\) & angle variable conjugate to \(I_j\); ordinary polar/Euler angles are locally defined & Chapters~\ref{chap:rigid-integrability}, \ref{chap:perturbation-chaos}\\
\(\mathbb T^n\) & \(n\)-torus, usually \((\R/2\pi\Z)^n\) & Chapters~\ref{chap:rigid-integrability}, \ref{chap:perturbation-chaos}\\
\(\varepsilon\) & small perturbation parameter & Chapter~\ref{chap:perturbation-chaos}\\
\(H_0,H_1\) & integrable Hamiltonian and perturbing Hamiltonian & Chapter~\ref{chap:perturbation-chaos}\\
\(\omega(I)\) & frequency vector in action-angle variables & Chapter~\ref{chap:perturbation-chaos}\\
\(k\) & integer Fourier/resonance vector in KAM sections; spring/gravity constants are locally defined & Chapter~\ref{chap:perturbation-chaos}\\
\(\alpha,\tau\) & Diophantine constants; \(\alpha\) also appears as a wave number in stability sections only locally & Chapter~\ref{chap:perturbation-chaos}\\
\(K\) & standard-map kick strength & Chapter~\ref{chap:perturbation-chaos}\\
\(\rho\) & mass density in fluids and continua; radial coordinate only when locally stated & Chapters~\ref{chap:fluid}, \ref{chap:elastic-gauge}\\
\(u(x,t)\) & Eulerian fluid velocity field & Chapter~\ref{chap:fluid}\\
\(p(x,t)\) & fluid pressure in fluid chapters; canonical momentum elsewhere & Chapter~\ref{chap:fluid}\\
\(\sigma\) & Cauchy stress tensor in continuum chapters & Chapters~\ref{chap:fluid}, \ref{chap:elastic-gauge}\\
\(\mu,\nu\) & dynamic and kinematic viscosity in fluids; reduced mass \(\mu\) in Kepler sections & Chapters~\ref{chap:fluid}, \ref{chap:perturbation-chaos}\\
\(\mathrm{Re}\) & Reynolds number \(UL/\nu\) & Chapter~\ref{chap:fluid}\\
\(B\) & reference body/material domain in elasticity; bending stiffness in rod sections when locally defined & Chapter~\ref{chap:elastic-gauge}\\
\(X^A\) & material coordinates on \(B\) & Chapter~\ref{chap:elastic-gauge}\\
\(x^i\) & spatial coordinates in \(\R^3\) & Chapter~\ref{chap:elastic-gauge}\\
\(\varphi(X,t)\) & deformation map from reference body to space & Chapters~\ref{chap:fluid}, \ref{chap:elastic-gauge}\\
\(F^i{}_A\) & deformation gradient \(\partial_A\varphi^i\) & Chapter~\ref{chap:elastic-gauge}\\
\(C_{AB},E_{AB}\) & right Cauchy-Green tensor and Green strain & Chapter~\ref{chap:elastic-gauge}\\
\(P^A{}_i\) & first Piola-Kirchhoff stress & Chapter~\ref{chap:elastic-gauge}\\
\(\mathcal S\) & shape space for deforming bodies & Chapter~\ref{chap:elastic-gauge}\\
\(\mathcal A\) & mechanical/gauge connection on shape space & Chapter~\ref{chap:elastic-gauge}\\
\(\mathcal F\) & curvature of the connection \(\mathcal A\) & Chapter~\ref{chap:elastic-gauge}\\
\end{longtable}
\normalsize

\section{Common Overloads And How To Read Them}
\label{sec:overloaded-symbols}

\begin{center}
\begin{tabular}{p{0.16\textwidth}p{0.76\textwidth}}
\toprule
Symbol & Reading rule\\
\midrule
\(\omega\) & In symplectic sections it is the symplectic form. In rigid body
sections the spatial angular velocity is written in prose as spatial angular
velocity or \(\omega_{\mathrm{space}}\) when ambiguity matters. In fluids it is
vorticity. In KAM sections \(\omega(I)\) is a frequency vector.\\
\(I\) & In Chapter~\ref{chap:notation} and Chapter~\ref{chap:lag-ham}, \(I\)
may be a time interval. In rigid body mechanics \(I\) is an inertia tensor or
principal moment. In action-angle theory \(I_j\) is an action variable.\\
\(J\) & In mechanics \(J\) is usually angular momentum or a momentum map. In
continuum mechanics \(J=\det F\) is the deformation Jacobian.\\
\(p\) & In Hamiltonian mechanics \(p_i\) is canonical momentum. In fluid
mechanics \(p(x,t)\) is pressure.\\
\(F\) & In integrability \(F_i\) are first integrals. In elasticity
\(F^i{}_A\) is the deformation gradient.\\
\(A\) & In Kepler theory \(A\) may denote the Laplace-Runge-Lenz vector. In
gauge theory \(A\) or \(\mathcal A\) denotes a connection.\\
\(\lambda\) & A Lagrange multiplier in constrained mechanics, a second
viscosity coefficient in fluids, or a Lame modulus in elasticity; the section
defines which one.\\
\bottomrule
\end{tabular}
\end{center}

\section{A Map Of The Course}
\label{sec:course-map}

The course moves through four enlargements of the same idea. The diagram is not
a dependency graph; it is a loop that the notes revisit at increasing levels of
structure. A variational principle produces equations, Hamiltonian geometry
organizes their phase space, perturbation theory tests the stability of that
organization, and continuum mechanics returns to the same questions with
fields instead of finitely many coordinates.

\begin{figure}[htbp]
\centering
\begin{tikzpicture}[
  node distance=1.1cm and 1.35cm,
  box/.style={draw, rounded corners=2pt, align=center, minimum width=3.3cm,
  minimum height=0.92cm, fill=blue!4},
  smallbox/.style={draw, rounded corners=2pt, align=center, minimum width=2.8cm,
  minimum height=0.75cm, fill=gray!8},
  arr/.style={-{Latex[length=2.5mm]}, thick},
  softarr/.style={-{Latex[length=2.2mm]}, thick, gray!70}
]
\node[box] (lag) {Lagrangian mechanics\\configuration, action, variation};
\node[box, right=of lag] (ham) {Hamiltonian mechanics\\phase space, symplectic form};
\node[box, below=of ham] (chaos) {Integrability and chaos\\tori, resonances, transport};
\node[box, left=of chaos] (cont) {Fluids and elasticity\\fields, stress, frames};
\node[smallbox, below=of lag, xshift=-0.35cm] (examples) {Kepler and rigid body\\first exact models};
\node[smallbox, above=of ham, xshift=0.35cm] (demos) {Demos and diagnostics\\drift, sections, probability};
\draw[arr] (lag) -- (ham);
\draw[arr] (ham) -- (chaos);
\draw[arr] (chaos) -- (cont);
\draw[arr] (cont) -- (lag);
\draw[softarr] (examples) -- (lag);
\draw[softarr] (examples) -- (chaos);
\draw[softarr] (demos) -- (ham);
\draw[softarr] (demos) -- (chaos);
\end{tikzpicture}
\caption{The course loop. The central theories are connected by examples and
computational diagnostics: exact models teach the geometry, while numerical
demos test which structures survive discretization or perturbation.}
\label{fig:course-loop}
\end{figure}

\section{Units And Dimensions}
\label{sec:units-dimensions}

Dimensional consistency is treated as a mathematical constraint. In celestial
mechanics demos we often use astronomical units, years, and solar masses, so
\[
  G = 4\pi^2
\]
when the central mass is measured in solar masses. In rigid body mechanics,
moments of inertia have dimensions \(ML^2\), angular velocity has dimension
\(T^{-1}\), angular momentum has dimension \(ML^2T^{-1}\), and kinetic energy
has dimension \(ML^2T^{-2}\).

\section{How To Read The Figures}
\label{sec:read-figures}

Figures in the notes are schematic unless explicitly generated from a numerical
demo. Their purpose is to identify variables, spaces, and geometric relations.
For quantitative plots, the corresponding Python or Mathematica script should be
treated as part of the figure caption: the model assumptions and numerical
parameters matter.

\section{How To Read Long Derivations}
\label{sec:read-long-derivations}

Long derivations in mechanics usually have two layers. The first layer is
formal: differentiate an action, apply a balance law, invert a Legendre
transform, or change coordinates. The second layer is structural: identify the
conserved quantity, the reduced phase space, the constraint force, or the
connection hidden in the algebra. The notes spell out both layers whenever the
second layer is easy to miss.

A useful way to read each derivation is to keep a small ledger:
\[
\begin{array}{c|c}
\text{input} & \text{output}\\
\hline
\text{model assumptions} & \text{equation of motion or constraint}\\
\text{symmetry} & \text{conserved quantity or reduced variable}\\
\text{constitutive law} & \text{closed continuum PDE}\\
\text{choice of frame} & \text{connection and curvature}
\end{array}
\]
Most mistakes in mechanics are ledger mistakes: using an equation outside the
assumptions that produced it, changing a frame without transforming the
components, or interpreting a reduced equation as if no reduction had occurred.

When a derivation contains a local coordinate calculation, the conclusion should
be read as coordinate-independent only after the geometric object has been
identified. For example, the symbols \(\Gamma^\ell{}_{ij}\), \(\lambda_a\),
\(A_A\), and \(\omega^a{}_b\) are coordinate or frame components; the geodesic
equation, constraint reaction, gauge connection, and torsion are the underlying
objects.
